\begin{comment}

TO CHECK!!!

The complete costs are detailed in Table~\ref{tab:equip}.

-	Les hores emprades per la realització del treball (les corresponents als crèdits del TFE). Millor 
si estan separades per les tasques executades segons la planificació. En una primera aproximació podeu considerar
 que una de les vostres hores es paga a 15 €/h.
-	Les despeses operatives: electricitat, calefacció, aigua o telefonia. Sumeu la part proporcional dels termes fixos.
 També tingueu en compte els viatges, les despeses d'oficina i d'altres.
-	Les despeses experimentals: materials, aigua, electricitat, combustibles, llicències, llibres.
-	Si és el cas, amortització dels equips. Considereu que un PC s'amortitza linealment en 5 anys i un telèfon mòbil en 3 anys.
-	Serveis abonats per obtenir algunes dades (per exemple d'un servei de microscòpia electrònica, encara que no ho hagi pagat l'estudiant).
-	I totes aquelles que haurien de constar si el treball fos realista i s'hagués de facturar. 


\begin{table}[!htb]\centering
	\caption{Total Costs.}
	\begin{tabular}{ccc|c}
		\hline
		\textbf{Concept} & \textbf{Unit cost (\texteuro/h)} & \textbf{Quantity (h)} & \textbf{Total (\texteuro)} \\
		\hline
		\hline
		Development engineer & 25.00 & 1,040 & 26,000.00 \\ 
		Supervisor & 30.00 & 260 & 7,800.00 \\
		\hline
		\textbf{Total} & & & 33,800.00 \\
		\hline
	\end{tabular}
	\label{tab:equip}
\end{table}
\end{comment}


This section provides an estimated budget for the development of the present master's thesis. 
While the work has been carried out as part of professional activity at eRoots Analytics, 
the following analysis reflects the costs that would reasonably be incurred if the project were to be 
considered as an engineering activity. The objective of this section is therefore 
to quantify the economic resources associated with the execution of the work, including labour, operational 
expenses, equipment usage, and other indirect costs.

The main cost associated with this project corresponds to the time devoted to its development. 
The work has been carried out during regular working hours, with a full-time dedication equivalent 
to eight hours per day during a total time of 22 weeks which is equivalent to 880~h. In line with the academic requirements of the Master's Thesis, the total workload 
is assumed to correspond to the standard number of credits associated with the TFE.

For the purpose of this estimation, the project guidelines suggest an hourly cost of 15~€/h. However, the real hourly cost
is about 23~€/h and therefore this second cost is used for the budget calculation.
The total effort is distributed among the main tasks defined in the project planning, including literature review, 
model development, implementation of RMS and EMT simulation tools, validation and testing, and report writing.

\begin{table}[H]
\centering
\caption{Estimated labour cost}
\begin{tabular}{l c c c}
\hline
Task category & Estimated hours & Unit cost [€/h] & Cost [€] \\
\hline
Literature review and theoretical study & 160 & 23 & 3680 \\
Software development and implementation & 400 & 23 & 9200 \\
Simulation, testing and validation & 200 & 23 & 4600 \\
Documentation and thesis writing & 120 & 23 & 2760 \\
\hline
Total labour cost & 880 & -- & 20240 \\
\hline
\end{tabular}
\end{table}

The execution of the project entails operational costs related to electricity consumption, HVAC usage, 
internet connectivity and general office resources. As the work was performed in a shared office 
environment with approximately 20 occupants, only the proportional share attributable to the thesis 
activity has been considered.

Office energy consumption values are derived from real data about the office consumption. eRoots pays around 100~€/month of electricity and
80€~/h of internet. 100~€/month of miscellaneous expenses are assumed.

\begin{table}[H]
\centering
\caption{Estimated operational expenses}
\renewcommand{\arraystretch}{1.15}
\small
\begin{tabular}{|l|c|}
\hline
\textbf{Concept} & \textbf{Estimated cost [€]} \\
\hline
Electricity and HVAC (proportional share) & 25 \\
Internet and telecommunication services & 20 \\
Office consumables and miscellaneous expenses & 25 \\
\hline
Total operational expenses & 70 \\
\hline
\end{tabular}
\end{table}

The thesis was developed using company-owned electronic equipment, namely a work laptop and an external 
monitor. The technical characteristics and typical power consumption of these devices are consistent 
with commercially available mid-range equipment, such as the ASUS Vivobook series and standard AOC 
monitors \cite{AsusVivabook,AOC}.

A linear amortisation model has been applied, assuming a useful life of five years for both devices. 
This assumption is consistent with common practice in life-cycle assessments and office energy studies 
\cite{ADEME2023Offices,EEA2021OfficeEnergy}. In line with the environmental impact assessment, an allocation 
factor of 0,084 has been used, corresponding to five months of use over a five-year lifespan.

\begin{table}[H]
\centering
\caption{Equipment amortization}
\renewcommand{\arraystretch}{1.15}
\small
\begin{tabular}{|l|c|c|c|}
\hline
\textbf{Equipment} & \textbf{Purchase value [€]} & \textbf{Allocation factor} & \textbf{Allocated cost [€]} \\
\hline
Laptop & 1000 & 0,084 & 84 \\
External monitor & 150 & 0,084 & 12,6 \\
\hline
Total equipment amortization & -- & -- & 96,6 \\
\hline
\end{tabular}
\end{table}

The project relies exclusively on numerical simulation and software development. All tools employed are 
either open-source or licensed by the hosting company, and no additional paid licenses, experimental 
services, or external data sources were required specifically for the execution of this work. As a result, 
no direct costs are allocated to software and external services.

\subsection*{Total estimated budget}

Table~\ref{tab:total_budget} summarizes the total estimated cost of the project, combining labour, 
operational expenses and equipment amortization. This budget represents an approximation of the 
resources required to carry out the thesis in a professional engineering context.

\begin{table}[H]
\centering
\caption{Total estimated project budget}
\label{tab:total_budget}
\renewcommand{\arraystretch}{1.15}
\small
\begin{tabular}{|l|c|}
\hline
\textbf{Cost category} & \textbf{Cost [€]} \\
\hline
Labour costs & 20240 \\
Operational expenses & 70 \\
Equipment amortisation & 96,6 \\
\hline
\textbf{Total estimated budget} & \textbf{20\,406,6} \\
\hline
\end{tabular}
\end{table}